\section*{Conclusion}
\addcontentsline{toc}{section}{Conclusion}

L'étude par éléments finis présentée dans ce document avait pour objectif de vérifier la
tenue mécanique du mécanisme de séparation de la fusée SP-02 avant l'usinage des pièces.\\

L'analyse statique de la pièce principale de séparation, réalisée sous une charge
représentative de la masse de la fusée au décollage affectée d'un facteur de sécurité de
$1,5$, conduit à une déformée maximale de $2,04$~mm. Les contraintes de von Mises restent
faibles sur la majeure partie de la pièce, mais atteignent localement environ $600$~MPa,
au-delà de la limite élastique de l'alliage 7075-T6. Cette valeur étant concentrée sur un
nombre restreint de nœuds, au voisinage immédiat des encastrements, elle relève très
probablement d'une singularité numérique liée à l'idéalisation des liaisons vissées plutôt
que d'un risque réel de plastification. Elle constitue néanmoins le point de vigilance
principal de cette étude et justifierait soit un raffinement local du modèle, soit une
modélisation plus fidèle des vis M5.\\

Les loquets présentent quant à eux des niveaux de sollicitation très faibles, avec des
déformées de l'ordre de $10^{-8}$~mm. La marge disponible sur ces pièces est donc
confortable, ce qui est cohérent avec leur géométrie massive au regard des efforts repris.
L'analyse modale conduite sur la pièce support, en configuration libre, a principalement
servi à prendre en main la résolution SOL-103 ; elle ne constitue pas encore une
vérification du comportement vibratoire du mécanisme en conditions de vol.\\

La confrontation aux calculs analytiques appelle une lecture nuancée. Sur la pièce de
séparation, l'estimation par la théorie des plaques circulaires encastrées donne
$0,26$~mm contre $2,04$~mm en simulation, soit un écart d'un facteur huit, très au-delà
du critère de $5~\%$. Cet écart n'invalide pas le modèle numérique : il traduit
l'inadéquation du modèle analytique, qui suppose une plaque d'épaisseur constante
encastrée sur tout son contour, là où la pièce réelle est une structure conique fixée en
quatre points et chargée de manière non uniforme. Le calcul analytique doit donc être
compris ici comme une vérification d'ordre de grandeur et non comme une validation
croisée. Sur les loquets, le calcul de poutre encastrée donne $0,45$~mm pour une longueur
de $50$~mm, alors que la longueur réellement en contact dans la géométrie usinée est de
$6$~mm, ce qui explique l'essentiel de l'écart observé avec la simulation.\\

Le mécanisme de séparation est donc jugé apte à reprendre les charges de vol, sous réserve
des points identifiés ci-dessus. Plusieurs travaux restent toutefois à mener pour consolider
cette conclusion : une étude de convergence en maillage documentée sur la pièce de
séparation, une analyse modale conduite avec les conditions aux limites réelles de
l'assemblage et confrontée au spectre d'excitation attendu en vol, la prise en compte des
efforts aérodynamiques en complément des seuls efforts d'inertie, et enfin une validation
expérimentale du mécanisme assemblé. Ces limites étant assumées, les résultats obtenus
fournissent une base suffisante pour engager la fabrication des pièces.
